\documentclass[12pt]{article}

\setlength{\parindent}{0pt}
\setcounter{secnumdepth}{3}
\usepackage[margin=1in]{geometry}

% packages
\usepackage{amsmath, amsfonts, amsthm, amssymb}
\usepackage{graphicx, float}
\usepackage{mathtools}
\usepackage{titlesec}
\usepackage{interval}
\usepackage{hyperref}
\usepackage{siunitx}
\usepackage{titling}
\usepackage{vwcol}
\usepackage{setspace}
\usepackage{empheq}
\usepackage{cancel}
\usepackage{esdiff}
\usepackage{multicol}
\usepackage{mdframed}
\usepackage{esdiff}
\usepackage{tikzsymbols}
\usepackage{multicol}
\usepackage{tikz}
\usepackage{varwidth}
\usepackage{pgfplots}
\usepackage{fancyhdr}
\usepackage{enumitem}
\usepackage{tcolorbox}

\intervalconfig {
	soft open fences
}

% header
\pagestyle{fancy}
\fancyhf{}
\lhead{Ethan Fan}
\rhead{AP Calculus BC}
\cfoot{\thepage}

% section format
\titleformat{\section}
  {\large \bfseries}
  {}
  {0em}
  {}
\titlespacing*{\section}{0pt}{0.5em}{0.3em}

\titleformat{\subsection}[runin]
  {\bfseries}
  {}
  {0em}
  {}
\titlespacing*{\subsection}{0pt}{0pt}{0em}

\titleformat{\subsubsection}[runin]
  {\itshape}
  {(\roman{subsubsection})}
  {0em}
  {}

% commands
\newcommand{\alignedintertext}[1]{%
  \noalign{%
    \vskip\belowdisplayshortskip
    \vtop{\hsize=\linewidth#1\par
    \expandafter}%
    \expandafter\prevdepth\the\prevdepth
  }%
}

\newcommand*{\paren}[1]{\ensuremath\left(#1\right)}
\newcommand*{\limit}[2][x]{\ensuremath{\displaystyle\lim_{#1\to#2}}}
\newcommand*{\Limit}[3][x]{\ensuremath{\displaystyle\lim_{#1\to#2}\left[#3\right]}}
\newcommand*{\deriv}[1][x]{\ensuremath{\dfrac{\mathrm{d}}{\mathrm{d}#1}}}
\newcommand*{\Deriv}[2][x]{\ensuremath{\dfrac{\mathrm{d}}{\mathrm{d}#1}\left[#2\right]}}
\newcommand{\ddx}{\dfrac{d}{dx}}
\newcommand{\dydx}{\dfrac{dy}{dx}}
\newcommand{\dydt}{\dfrac{dy}{dt}}
\newcommand{\dxdt}{\dfrac{dx}{dt}}
\newcommand*{\iinteg}[2][x]{\ensuremath{\displaystyle\int #2\;\mathrm{d}#1}}
\newcommand*{\dinteg}[4][x]{\ensuremath{\displaystyle\int_{#2}^{#3}#4\;\mathrm{d}#1}}
\newcommand*{\abs}[1]{\ensuremath{\left|#1\right|}}
\newcommand*{\eps}{\varepsilon}
\newcommand*{\real}{\ensuremath\mathbb{R}}
\newcommand*{\integer}{\ensuremath\mathbb{Z}}
\newcommand*{\posint}{\ensuremath\mathbb{Z^+}}
\newcommand*{\cbrt}[1]{\ensuremath{\sqrt[3]{#1}}}
\newcommand*{\lheqzero}{\ensuremath{\underset{\text{L'H}}{\overset{\left[\frac00\right]}{=}}}}
\newcommand*{\lheqinfty}{\ensuremath{\underset{\text{L'H}}{\overset{\left[\frac{\infty}{\infty}\right]}{=}}}}
\newcommand{\tor}{\text{ or }}
\newcommand{\floor}[1]{\lfloor #1 \rfloor}

\newcommand{\vem}{\vspace{0.5em}}
\newcommand{\example}[1]{\section{Example #1}}
\newcommand{\problem}[1]{\section{Problem #1}}
\newcommand{\subproblem}[1]{\subsection{(#1)} \,}

\DeclareMathOperator{\dne}{DNE}
\DeclareMathOperator{\sgn}{sgn}

\newcommand{\definition}[1]{\begin{tcolorbox}[colback=red!5!white,colframe=red!75!black,parbox=false] #1 \end{tcolorbox}}

\newcommand{\theorem}[2]{\begin{tcolorbox}[title={#1},colback=blue!5!white,colframe=blue!75!black,parbox=false] #2 \end{tcolorbox}}

\allowdisplaybreaks
% \postdisplaypenalty=100000

% document
\begin{document}

\vspace*{0cm}

% opening
\begin{center}
    {\Large \bfseries Problem Set 6} \\[0.5cm]
    {\large The Fundamental Theorem of Calculus} \\[0.35cm]
    {\normalsize September 17, 2026}
\end{center}

% problems

\problem{3}

\subproblem{c}
Both formulas remain the same. \vem

\problem{4}

\subproblem{b}
\begin{gather*}
    \deriv  \int_1^{x^4} \sec t\, dt\\
    =\boxed{4x^3\sec x^4}
\end{gather*}

\subproblem{c}
\begin{gather*}
    \deriv \int_x^\pi{ \sqrt{1+\sec t}} \, dt\\
    = \boxed{- \sqrt{1+\sec x}}
\end{gather*}

\subproblem{e}
\begin{gather*}
    \deriv \int_1^{\tan x} \sqrt{t+\sqrt{t}} \, dt\\
    =\boxed{\sec ^2 x \sqrt{\tan x+\sqrt{\tan x}}}
\end{gather*}

\subproblem{f}
\begin{gather*}
    \deriv \int_{1-3x}^1 \frac{v}{v^3+1} \, dv\\
    =\boxed{\frac{3(1-3x)}{(1-3x)^3+1}}
\end{gather*}

\problem{6}

\subproblem{a}
\begin{gather*}
    F'(x)=\ddx \int_1^x f(t)\, =f(x)\\
    F''(x)=f'(x)= \ddx \int_1^{x^2} \frac{\sqrt{1+u^4}}{u} \, du =2x \frac{\sqrt{1+x^8}}{x^2}\\
    F''(2)=2\cdot 2 \cdot \frac{\sqrt{1+2^8}}{4}=\boxed{\sqrt{257}}
\end{gather*}

\subproblem{b}
\begin{gather*}
    y'=\ddx \int_0^x \frac{1}{1+t+t^2} \, dt= \frac{1}{1+x+x^2}\\
    y''=\ddx \frac{1}{1+x+x^2}=-\frac{2x+1}{(1+x+x^2)^2}\\
    y''>0 \implies 1+2x<0 \implies \boxed{x\in (-\infty , -\frac{1}{2})}
\end{gather*}

\subproblem{c}
\begin{gather*}
    6+\int_a^x \frac{f(t)}{t^2} \, dt=2\sqrt{x}\\
    \ddx \paren{6+\int_a^x \frac{f(t)}{t^2} \, dt} = \ddx 2\sqrt{x}\\
    \frac{f(x)}{x^2}=\frac{1}{\sqrt{x}}\\
    \boxed{f(x)=x\sqrt{x}}\\
    6+\int_a^x \frac{1}{\sqrt{x}} \, dt=2\sqrt{x}\\
    x=a \implies 6+0 = 2\sqrt{a}\\
    \boxed{a=9}
\end{gather*}

\problem{7}

\subproblem{a}
\begin{gather*}
    \int_1^2 h''(u) \, du= h'(2)-h'(1)=5-2=3
\end{gather*}
The integral of $h''(u)$ only depends on its anti-derivative, which is $h'(u)$.\vem

\subproblem{b}
\begin{gather*}
    \int_1^2 h'(u) \, du= h(2)-h(1)=6-(-2)=8
\end{gather*}

\subproblem{c}
\[
V(20) - V(0) = \int_0^{20} V'(t)\, dt = \int_0^{20} r(t)\, dt
\]
The initial value does not matter, as the answer only measures the net change in volume.

\problem{8}

\subproblem{a}
As the function is positive for all values of $x\ne 0$, the area under the curve cannot possibly be negative. It fails because the function is not continuous on the interval.\vem

\subproblem{b}

\vem
(i) 
\begin{gather*}
    \frac{g(x+h)-g(x)}{h}=\frac{1}{h}\paren{\int _a^{x+h} f(t) \, dt-\int _a^{x} f(t) \, dt}=\frac{1}{h} \int _x^{x+h} f(t) \, dt
\end{gather*}

(ii)
\begin{gather*}
    m_h \le f(t) \le M_h\\
    \int _x^{x+h} m_h \, dt\le \int _x^{x+h} f(t) \, dt \le \int _x^{x+h} M_h \, dt\\
    h \cdot m_h \le g(x+h)-g(x)\le h \cdot M_h\\
    m_h\le \frac{g(x+h)-g(x)}{h} \le M_h
\end{gather*}

(iii)
\begin{gather*}
    \lim_{h \to 0^+} m_h=\lim_{h \to 0^+}M_h = f(x)\\
    \lim_{h \to 0^+} g'(x)=\lim_{h \to 0^+}\frac{g(x+h)-g(x)}{h}=f(x)
\end{gather*}
For $h<0$, both the interval of integration and the value of $h$ changes, but the result remains the same. \vem

\subproblem{c}
\begin{gather*}
    F'(x)=g'(x)=f(x)\\
    F(x)=g(x)+C\\
    F(a)=g(a)+C=\int_a^a f(t)\, dt+C=C\\
    F(b)=g(b)+C=\int_a^b f(t)\, dt+C\\
    \int_a^b f(t)\, dt=F(b)-C=F(b)-F(a)
\end{gather*}

\subproblem{d}

\vem
(i)
\begin{gather*}
    x \in [-1,0] \implies g(x)=0\\
    x\in (0,1] \implies g(x)=\int_0^x 1\, dt=x\\
    g(x)=\begin{cases}
        0,\quad x\in [-1,0]\\
        x, \quad x\in (0,1]
    \end{cases}
\end{gather*}

(ii)
\begin{gather*}
    g(0)=\lim_{x\to 0^-} g(x)=\lim_{x\to 0^+} g(x)=0\\
    \lim_{x\to 0^-} g'(x)=0\\
    \lim_{x\to 0^+} g'(x)=1\\
    \lim_{x\to 0^-} g'(x) \ne \lim_{x\to 0^+} g'(x)
\end{gather*}

(iii)
As $f$ is not continuous on the interval, differentiability is not guaranteed for $g(x)$. The accumulation is good, as it renders $g(x)$ a continuous function. However, sharp corners in the graph still exist.















\end{document}

